% INR icon (v2) -- shared learned-row grid (code x=-0.90 | proc x=-0.08 | liver x=0.74).
% Identity: COORDINATE-BASED CONTINUOUS FUNCTION. Input = spatial query (x,y,z) as a 3-cell stack
% (same footprint as b/z); processor = small network f_theta (its WEIGHTS are the shape); output =
% a continuous liver field. Distinct from latent: fed a coordinate, outputs a field.
% Consumes \liveroutline, \liverfield (learned_common).
\newcommand{\iconINR}{%
  % coordinate query (x,y,z) -- 3 cells matching the other code stacks
  \foreach \i/\l in {0/x, 1/y, 2/z}{
    \draw[fill=white] (-0.99,{0.10-\i*0.165}) rectangle ++(0.18,0.13);
    \node[font=\tiny] at (-0.90,{0.165-\i*0.165}) {$\l$};
  }
  \node[font=\tiny, anchor=south] at (-0.90,0.30) {$\mathbf{x}$};
  \draw[-{Latex[length=1mm]}] (-0.78,0) -- (-0.40,0);
  % small MLP f_theta (edges first, then nodes), centred at x=-0.08
  \foreach \hy in {0.17,0,-0.17}{
    \draw[gray, very thin] (-0.28,0.10) -- (-0.07,\hy);
    \draw[gray, very thin] (-0.28,-0.10) -- (-0.07,\hy);
    \draw[gray, very thin] (-0.07,\hy) -- (0.14,0);
  }
  \foreach \p in {(-0.28,0.10),(-0.28,-0.10),(-0.07,0.17),(-0.07,0),(-0.07,-0.17),(0.14,0)}{
    \node[circle, draw, inner sep=0.5pt, scale=0.5, fill=white] at \p {};
  }
  \node[font=\tiny, anchor=south] at (-0.08,0.30) {$f_\theta$};
  \draw[-{Latex[length=1mm]}] (0.20,0) -- (0.45,0);
  % output liver, EVALUATED at one specific query point p=(x,y,z)
  % (flat glyph -- same as icons 6/7/9; no radial-gradient "shine")
  \begin{scope}[shift={(0.74,0)}, scale=0.45] \liverglyph \end{scope}
  \coordinate (qp) at (0.80,-0.03);
  % crosshair query-point marker
  \draw[black, thin] ($(qp)+(-0.06,0)$) -- ($(qp)+(0.06,0)$);
  \draw[black, thin] ($(qp)+(0,-0.06)$) -- ($(qp)+(0,0.06)$);
  \fill[black] (qp) circle (0.021);
  \draw[white, line width=0.25pt] (qp) circle (0.021);
}
